\clearpage\thispagestyle{empty}
\pdfbookmark[0]{Quellen zu Gauß’ Meridianbeobachtungen der Iris (September 1847)}{iris613-sources}
\begin{center}\textbf{Quellen zu Gauß’ Meridianbeobachtungen der Iris (September 1847)}\end{center}
\noindent Diese KI-Übertragung ins heutige Deutsch folgt dem Auszug auf S.~474 der \href{https://gdz.sub.uni-goettingen.de/id/PPN236007467}{\textit{Werke}, Band~VI} (Göttingen, 1874). Der Text der Werkausgabe beruht auf einem Brief von Carl Friedrich Gauß an Heinrich Christian Schumacher vom 5. September 1847 (\href{https://api.digitale-sammlungen.de/iiif/image/v2/bsb10049451_00105/full/1800,/0/default.jpg}{Sp.~193--194}), einem Bericht über spätere Beobachtungen (\href{https://api.digitale-sammlungen.de/iiif/image/v2/bsb10049451_00107/full/1800,/0/default.jpg}{Sp.~197--198}) und einer weiteren Beobachtung in den vermischten Nachrichten (\href{https://api.digitale-sammlungen.de/iiif/image/v2/bsb10049451_00112/full/1800,/0/default.jpg}{Sp.~207--208}). Alle drei stehen in den \textit{Astronomischen Nachrichten}, Band~XXVI, Nr.~613, laut Schlusszeile auf Sp.~208 ausgegeben am 14. Oktober 1847.
\par\smallskip\noindent Die Zeitschrift druckt im Brief die Elemente von Gould und Goldschmidt getrennt ab. Der Auszug der Werkausgabe setzt an diese Stelle vier Striche. Die Zahlenschreibweise, die fünf Beobachtungszeilen und die Einschränkungen zur Sichtbarkeit folgen hier der Werkausgabe. Die Lesart \textit{Din Luft} ist beibehalten. KI-Übertragung und KI-Satz.
\clearpage
